\documentclass{../base/canon}

%% canon_variables.tex — Valores por defecto de metadatos institucionales
%%
%% Incluir ANTES del \begin{document} para sobreescribir los defaults de la clase.
%% El template puede sobreescribir cualquier valor después de este \input.
%%
%% Uso en template:
%%   %% canon_variables.tex — Valores por defecto de metadatos institucionales
%%
%% Incluir ANTES del \begin{document} para sobreescribir los defaults de la clase.
%% El template puede sobreescribir cualquier valor después de este \input.
%%
%% Uso en template:
%%   %% canon_variables.tex — Valores por defecto de metadatos institucionales
%%
%% Incluir ANTES del \begin{document} para sobreescribir los defaults de la clase.
%% El template puede sobreescribir cualquier valor después de este \input.
%%
%% Uso en template:
%%   \input{../base/canon_variables}
%%   \SetDocTitle{Mi Informe Específico}   % sobreescribe el default

\SetDocLogo{../assets/logo.pdf}
\SetDocUnit{Tecnología \& Sistemas}
\SetContactUnit{Mesa de soporte}
\SetContactMembers{%
  Equipo CoreX & Soporte & soporte@empresa.com%
}
\SetContactNote{Para soporte técnico o consultas sobre este documento.}

%%   \SetDocTitle{Mi Informe Específico}   % sobreescribe el default

\SetDocLogo{../assets/logo.pdf}
\SetDocUnit{Tecnología \& Sistemas}
\SetContactUnit{Mesa de soporte}
\SetContactMembers{%
  Equipo CoreX & Soporte & soporte@empresa.com%
}
\SetContactNote{Para soporte técnico o consultas sobre este documento.}

%%   \SetDocTitle{Mi Informe Específico}   % sobreescribe el default

\SetDocLogo{../assets/logo.pdf}
\SetDocUnit{Tecnología \& Sistemas}
\SetContactUnit{Mesa de soporte}
\SetContactMembers{%
  Equipo CoreX & Soporte & soporte@empresa.com%
}
\SetContactNote{Para soporte técnico o consultas sobre este documento.}

\SetDocType{REPORTE TÉCNICO}
\SetDocTitle{Título del Reporte Técnico}
\SetDocCode{REP-TEC-000}
\SetDocArea{Tecnología \& Sistemas}
\SetDocAuthor{Autor Técnico}
\SetDocDate{\today}

\begin{document}

\section{Contexto técnico}

\begin{ParrafoMetodologico}
Describir el sistema, componente o proceso analizado. Incluir versiones, entorno
(producción/staging), período de observación y herramientas utilizadas.
\end{ParrafoMetodologico}

\section{Hallazgos}

\KeyFindingBox[Hallazgo principal]{
  Descripción concisa del hallazgo técnico más relevante.
}

\begin{itemize}
  \item Hallazgo técnico 1.
  \item Hallazgo técnico 2.
\end{itemize}

\WarningBox[Alerta crítica]{
  Describir aquí cualquier condición que requiera acción inmediata.
}

\section{Evidencia técnica}

\subsection{Consulta / Comando ejecutado}

\begin{CodeBlock}
-- Consulta o comando que generó la evidencia
SELECT * FROM tabla WHERE condicion = true;
\end{CodeBlock}

Ruta del endpoint o archivo: \CodePath{/ruta/del/recurso}

\subsection{Resultado}

\begin{table}[H]
  \centering
  \caption{Resultado de la verificación}
  \begin{tabular}{l l r}
    \toprule
    \textbf{Elemento} & \textbf{Estado} & \textbf{Valor} \\
    \midrule
    Elemento A & OK    & 0 \\
    Elemento B & Error & 0 \\
    \bottomrule
  \end{tabular}
\end{table}

\section{Recomendaciones técnicas}

\begin{enumerate}
  \item Acción correctiva 1 — responsable, fecha estimada.
  \item Acción correctiva 2 — responsable, fecha estimada.
\end{enumerate}

\SignatureBlock{Nombre Apellido}{Ingeniero de Sistemas}

\ContactSignature

\end{document}
